% \nv{...}: một ý trong cột Nhiệm vụ, dòng xuống hàng thụt vào thẳng với chữ đầu tiên.
\providecommand{\nv}[1]{\par\noindent\hangindent=1em\hangafter=1\makebox[1em][l]{--}#1}

\begin{table}[H]
\centering
\renewcommand{\tabularxcolumn}[1]{m{#1}}
\begin{tabularx}{\textwidth}{|c|c|>{\raggedright\arraybackslash}X|>{\centering\arraybackslash}m{1.9cm}|}
\hline
\textbf{MSSV} & \textbf{Họ và tên} & \centering\textbf{Nhiệm vụ} & \textbf{Mức độ đóng góp} \\ \hline
24120048 & Lê Hoàng Hiếu &
\nv{Tìm hiểu và báo cáo: so sánh cách dùng file descriptor trong \texttt{open()}, \texttt{dup()}, \texttt{fork()}; trace system call \texttt{fork}.}
\nv{Viết chương trình kiểm thử \texttt{user/sysinfotest.c}.}
\nv{Dựng khung báo cáo, viết mục thông tin nhóm, mức độ hoàn thành và kết quả kiểm thử.}
& 100\% \\ \hline
24120090 & Đặng Hồng Minh &
\nv{Tìm hiểu và báo cáo: các trường \texttt{pid}, \texttt{state}, \texttt{context} của \texttt{struct proc}; cơ chế lập lịch trong \texttt{scheduler()}, \texttt{sched()}, \texttt{yield()}.}
\nv{Đăng ký system call \texttt{sysinfo} (\texttt{syscall.h}, \texttt{syscall.c}, \texttt{usys.pl}, \texttt{user.h}, \texttt{Makefile}).}
\nv{Thêm bộ đếm \texttt{nsched} vào \texttt{struct proc}, viết \texttt{proccollect()} thu thập danh sách process.}
& 100\% \\ \hline
24120405 & Đinh Tiến Phát &
\nv{Tìm hiểu và báo cáo: trường \texttt{ofile[]} của \texttt{struct proc} và \texttt{file[]} của \texttt{ftable}; trace system call \texttt{open}.}
\nv{Thiết kế \texttt{struct sysinfo} (\texttt{kernel/sysinfo.h}), viết \texttt{sys\_sysinfo()}.}
\nv{Viết \texttt{filecollect()}, \texttt{fileidx()}, \texttt{procs\_using\_file()} thu thập danh sách open file.}
& 100\% \\ \hline
\end{tabularx}
\end{table}

Mỗi thành viên viết phần báo cáo tương ứng với phần mình tìm hiểu và cài đặt. Các mục 4.4, 5, 7, 8 được ghép từ phần đóng góp của cả ba người.
